%%%%%%%%%%  scan idx 35  |  folio 22  %%%%%%%%%%
\begin{equation}\tag{1.4}
\lambda_t(1)=1+t.
\end{equation}
On a alors
\begin{equation}\tag{1.5}
\lambda_t(n)=(1+t)^n,
\end{equation}
d'où $\lambda^i(n)=\binom{n}{i}$.

La
notion de $\lambda$-homomorphisme de $\lambda$-anneaux est claire~; un \emph{$\lambda$-anneau augmenté} est un $\lambda$-anneau muni d'un $\lambda$-homomorphisme dans $\Z$.

Pour définir la notion de $\lambda$-\emph{anneau spécial}, les considérations suivantes sont utiles. Soit $k$ un anneau commutatif avec unité, soit $K=1+k[[t]]^+$ le groupe des séries formelles à coefficients dans $k$, d'augmentation $1$~; on va y introduire une loi de $\lambda$-anneau dont la structure additive soit la multiplication des séries formelles. La multiplication de cet anneau sera notée $f\circ g$. Elle est entièrement déterminée par la condition que les coefficients $c_n=c_n(f\circ g)$ de $f\circ g$ soient donnés par des polynômes \emph{universels} à coefficients entiers en les coefficients $a_i=c_i(f)$ et $b_i=c_i(g)$ de $f$ et $g$ respectivement~:
\begin{equation}\tag{1.6}
c_n=P_n(a_1,\ldots,a_n;b_1,\ldots,b_n),
\end{equation}
qu'elle soit distributive par rapport à l'«~addition~» $fg$, et soit donnée pour des facteurs linéaires par~:
\begin{equation}\tag{1.7}
(1+at)\circ(1+bt)=1+abt.
\end{equation}
Les $P_n$ se calculent par des calculs de fonctions symétriques~; $P_n$ est isobare de poids $n$ en les $a_i$ et isobare de poids $n$ en les $b_i$ ($a_i$ et $b_i$ étant de poids $i$) et d'ailleurs symétrique en $a=(a_i)$ et $b=(b_i)$. De même, les $\lambda^i(f)$ sont déterminés sans ambiguïté par la condition que 
